% ============================================================================
% Conclusion
% ============================================================================

We presented \textbf{DUNE} (Dual-tag UWB Navigation Engine), an ultra-wideband
localization system delivering full 6-DOF pose estimation for a 13-DOF
rocker--bogie mobile manipulator operating in unstructured terrain.
DUNE addresses the fundamental limitations of all prior open-source Decawave
DW1000 implementations through three focused contributions.

\textbf{Dual-tag 6-DOF pose estimation.}
Two UWB tags at a \SI{50}{\centi\metre} baseline on the rover chassis provide
geometric yaw from their world-frame position difference.
Combined with five-anchor 3D ranging (one elevated anchor making $z$
observable) and BNO085 IMU fusion for roll and pitch, the system outputs
complete $(x, y, z, \phi, \theta, \psi)$ at \SI{\approx10}{\hertz}---the
first UWB implementation to provide full 6-DOF pose without motion capture.

\textbf{Two-tier in-situ calibration.}
Tier~1 antenna-delay binary search eliminates the 30--50\,cm bulk bias
present on uncalibrated boards, requiring only a tape measure and taking
under two minutes per board.
Tier~2 spatial error mapping corrects residual position-dependent multipath.
We document and remove a double-correction failure in naive prior approaches
in which the same bias $b_0$ was both pushed to NVS and subtracted in
software simultaneously, producing a net over-correction of $2b_0$ that
inverted the bias direction entirely.

\textbf{Same-frame NLOS detection and weighted multilateration.}
Reading first-path and receive power from the same RANGE\_REPORT frame on
the tag side yields a physically valid per-anchor NLOS gap signal, correcting
an error in all existing libraries where the gap was computed across two
different radio links.
Continuous weighting in the Levenberg--Marquardt solver, combined with an
independent MAD gate for gross outliers, delivers robust positioning under
partial obstructions.

\todo{Insert quantitative summary once hardware experiments complete:
DUNE achieves X cm 2D RMSE and Y deg heading RMSE at Z Hz on dynamic
rover trajectories over sandy terrain, compared to W cm for an
uncalibrated single-tag baseline.}

\textbf{Future work.}
Immediate next steps include extending the EKF from 2D to a 9-state
3D formulation once the five-anchor geometry is field-validated,
tighter IMU--UWB coupling directly on raw ranges, and coordinated
TDMA slot assignment for deterministic multi-tag operation.
Longer-term, migration to the Qorvo DW3000 (pin-compatible, production-current)
will extend the system lifetime beyond the DW1000 end-of-life.
